\documentclass[journal,twoside,web]{ieeecolor}
\usepackage{generic}

% Keep the current XeLaTeX-compatible setup to preserve template stability.
\usepackage[scheme=plain]{ctex}

\usepackage{cite}
\usepackage{amsmath,amssymb,amsfonts}
\usepackage{algorithmic}
\usepackage{graphicx}
\usepackage{textcomp}
\usepackage{mathrsfs}
\usepackage{booktabs}
\usepackage{algorithm}
\usepackage{booktabs}
\usepackage{makecell}
\usepackage{multirow}
\usepackage{graphicx}   % Required for figures
\usepackage{subcaption} % Required for subfigures
\usepackage{array}
%\usepackage{epstopdf}  % Optional: enable if pdflatex must convert EPS files
\renewcommand{\algorithmicrequire}{\textbf{Input:}}
\renewcommand{\algorithmicensure}{\textbf{Output:}}

\def\BibTeX{{\rm B\kern-.05em{\sc i\kern-.025em b}\kern-.08em
    T\kern-.1667em\lower.7ex\hbox{E}\kern-.125emX}}
\markboth{\journalname, VOL. XX, NO. XX, XXXX}
{Li \MakeLowercase{\textit{et al.}}: Title}
\begin{document}
\title{Preparation of Papers for IEEE Trans on Industrial Informatics (February 2022)1}
\author{Yulin Li, Chunyan Zhang,~\IEEEmembership{Member,~IEEE}, Qiuyang Fang, Xudong Zhao,~\IEEEmembership{Senior Member,~IEEE}, and Jianlei Zhang,~\IEEEmembership{Member,~IEEE}%
\thanks{Yulin Li, Chunyan Zhang, Qiuyang Fang, and Jianlei Zhang are with the College of Artificial Intelligence, Nankai University, Tianjin 300381, China (e-mail: yulinlifight@mail.nankai.edu.cn; zhcy@nankai.edu.cn; qiuyangfang@mail.nankai.edu.cn; jianleizhang@nankai.edu.cn).}%
\thanks{Xudong Zhao is with the Key Laboratory of Intelligent Control and Optimization for Industrial Equipment, Ministry of Education, Dalian University of Technology, Dalian 116024, China (e-mail: xdzhaohit@gmail.com).}%
\thanks{Corresponding author: Jianlei Zhang (e-mail: jianleizhang@nankai.edu.cn).}}

\maketitle

\begin{abstract}

phrases as this can result in a page being rejected by search engines.
Ensure that your abstract reads well and is grammatically correct.

\end{abstract}

\begin{IEEEkeywords}
Enter key words or phrases in alphabetical
order, separated by commas. For a list of suggested keywords, send a blank
e-mail to keywords@ieee.org or visit \underline
{http://www.ieee.org/organizations/pubs/ani\_prod/keywrd98.txt}
\end{IEEEkeywords}

\section{Introduction}
\label{sec:introduction}

\IEEEPARstart{T}{his}

\section{SYSTEM MODEL AND PROBLEM FORMULATION}

\begin{figure*}[!t]
 \centering
\includegraphics[width=\textwidth]{Main figure.eps}
 \caption{Overall framework and workflow of the proposed BHTA.}
 \label{fig:bhta_framework}
\end{figure*}

\begin{table}[htbp]
    \centering
    \caption{Summary of Main Notations}
    \label{tab:notations}
    \footnotesize
    \setlength{\tabcolsep}{4pt}
    \begin{tabular}{l p{4.8cm}}
        \toprule
        \textbf{Notation} & \textbf{Description} \\
        \midrule
        $N, M$ & Number of robots and tasks \\
        $Z, K$ & Number of capability dimensions and task types \\
        $\mathcal{R}$ & Set of robots \\
        $\mathcal{T}$ & Set of tasks \\
        $\mathcal{Z}$ & Set of capability dimensions \\
        $\mathcal{K}$ & Set of task types \\
        $\mathcal{P}$ & Task-priority list \\
        $\mathcal{L}_m$ & Demand vector of task $m$ \\
        $\mathcal{D}_k$ & Demand vector of task type $k$ \\
        $\mathcal{G}_n$ & Capability vector of robot $n$ \\
        $\mathcal{B}_{n}^m$ & Belief vector of robot $n$ over type of task $m$ \\
        $\widehat{\mathcal{V}}_{n,m}$ & Expected value of task $m$ by robot $n$ \\
        $\widehat{\mathcal{L}}_{n,m}$ & Expected demand vector of task $m$ by robot $n$ \\
        $\mathcal{SC}$ & Global overlapping coalition structure \\
        $\mathcal{A}_m$ & Coalition for task $m$ \\
        $A_m^{(n)}$ & Capability vector contributed by robot $n$ to task $m$ \\
        $\operatorname{Mem}(\mathcal{A}_m)$ & Member set of coalition for task $m$ \\
        $\|A_m^{(n)}\|_1$ & 1-norm of robot $n$'s allocation to task $m$ \\
        $\varsigma_{n,m}$ & Task completion degree estimated by robot $n$ \\
        $\mathcal{V}_k$ & Base value of task type $k$ \\
        $U$ & Global total utility \\
        $u_n^m(\mathcal{A}_m)$ & Expected utility of robot $n$ in coalition for task $m$ \\
        $E_{n,m}^{Cost}$ & Total cost of robot $n$ for task $m$ \\
        $E_{n,m}^{move}$ & Movement cost of robot $n$ to task $m$ \\
        $E_{n,m}^{wait}$ & Waiting cost of robot $n$ at task $m$ \\
        $E_{n,m}^{exec}$ & Execution cost of robot $n$ for task $m$ \\
        $\varpi, \gamma, \beta$ & Unit cost coefficients (move, wait, exec) \\
        $E_n^{max}$ & Energy budget of robot $n$ \\
        $AT_{n,m}$ & Arrival time of robot $n$ at task $m$ \\
        $ST_m$ & Synchronized start time of coalition for task $m$ \\
        $FT_{n,m}$ & Finish time of robot $n$'s part in task $m$ \\
        $FT_m$ & Finish time of coalition for task $m$ \\
        $T_{n,m}^{\mathrm{exec}}$ & Execution duration of robot $n$ for task $m$ \\
        $T_{n,m}^{\mathrm{pre}}$ & Pre-sync waiting time of robot $n$ at task $m$ \\
        $T_{n,m}^{\mathrm{post}}$ & Post-sync waiting time of robot $n$ at task $m$ \\
        \bottomrule
    \end{tabular}
\end{table}
\subsection{System Model}

Consider a group of $N$ heterogeneous robots, denoted by $\mathcal{R} = \{1, \dots, N\}$, deployed in an unknown environment to perform collaborative information collection over $M$ task locations $\mathcal{T} = \{1, \dots, M\}$.

\noindent\textit{1) Task Characterization:} Suppose that the environment contains $K$ task types, denoted by $\mathcal{K} = \{1, \dots, K\}$, and let $\mathcal{Z} = \{1,\ldots,z,\ldots,Z\}$ be the set of required capability dimensions. Each task type $k \in \mathcal{K}$ is associated with a potential demand vector for observation or service capabilities, denoted by $\mathcal{D}_k = [d_k^{(1)}, \dots, d_k^{(Z)}]$, and a base value $\mathcal{V}_k$. For task $m \in \mathcal{T}$, its specific demand vector is denoted by $\mathcal{L}_m = [l_m^{(1)}, \dots, l_m^{(Z)}]$. In addition, tasks are associated with a priority list $\mathcal{P} = \{\rho_1, \dots, \rho_M\}$, where $\rho_m$ denotes the priority weight of task $m$, and a larger value indicates a higher priority.

\noindent\textit{2) Robot Capability Characterization:} Each robot possesses reusable information-collection service capabilities, including onboard sensors, communication modules, and related functional units. For robot $n \in \mathcal{R}$, its capability vector is denoted by $\mathcal{G}_n = [\eta_n^{(1)}, \dots, \eta_n^{(Z)}]$, where $\eta_n^{(z)} \ge 0$ represents the maximum service capability that robot $n$ can provide in the $z$th dimension. If the robot does not possess that capability, then $\eta_n^{(z)} = 0$. To remain consistent with the subsequent mathematical notation, this paper still uses the term ``resource'' vector, but its physical meaning should be understood as reusable sensing and service capability rather than one-shot consumable inventory.
\subsubsection{Overlapping Coalition Model}
Since a single robot is often unable to satisfy the information-collection requirement of a task independently, multiple robots typically need to form a cooperative observation coalition around the task location. When the same robot supports multiple tasks, an overlapping coalition structure arises. We group the robots executing the same task into one coalition. For task $m$, the corresponding coalition is denoted by $\mathcal{A}_m = \{A_m^{(1)}, \dots, A_m^{(N)}\}$, where $A_m^{(n)} = \left[a_{n,m}^{(1)}, \dots, a_{n,m}^{(Z)}\right]$ represents the capability-allocation vector contributed by robot $n$ to task $m$. If robot $n$ does not participate in task $m$, then $A_m^{(n)} = \mathbf{0}$. Accordingly, the member set of task $m$ can be defined as
\begin{equation}
\operatorname{Mem}(\mathcal{A}_m) = \left\{ n \in \mathcal{R} \mid A_m^{(n)} \neq \mathbf{0} \right\}
\label{eq:coalition_members}
\end{equation}

\subsubsection{Uncertainty and Belief Modeling}

In practical environments, the robot swarm cannot know the true task demand $\mathcal{L}_m$ in advance. Moreover, due to sensor noise and observation-range limitations, individual robots may form different initial judgments about the type of the same task. Therefore, this paper introduces beliefs, i.e., probability distributions over task types, to characterize each robot's local cognition of task types and to estimate task values and capability demands. The local belief vector of robot $n$ for task $m$ is defined as $\mathcal{B}_{n}^m = [b_{n,1}^m, \dots, b_{n,K}^m]$, subject to the normalization condition $\sum_{k \in \mathcal{K}} b_{n,k}^m = 1$.
Based on this belief distribution, the expected base value of task $m$ estimated by robot $n$ is
\begin{align}
\widehat{\mathcal{V}}_{n,m} = \sum_{k \in \mathcal{K}} b_{n,k}^m \cdot \mathcal{V}_k,
\end{align}
and the expected demand estimate for task $m$ is
\begin{align}
\widehat{\mathcal{L}}_{n,m} = [\hat{l}_{n,m}^{(1)}, \dots, \hat{l}_{n,m}^{(Z)}].
\end{align}
For the $z$th demand dimension, the expectation is obtained by summing the task-type demand with local-belief weights:
\begin{align}
\hat{l}_{n,m}^{(z)} = \sum_{k \in \mathcal{K}} b_{n,k}^m \cdot d_k^{(z)}, \quad z \in \mathcal{Z}.
\end{align}

\subsection{Utility Function and Cost Model}
Individual robots make decisions by evaluating task utility. To achieve this, we first define the task completion degree. Under a proportional allocation rule, the revenue obtained by each robot is determined by its relative resource contribution to the current task coalition.


For robot $n$, the completion degree of task $m$ is evaluated as the average satisfaction across all effective resource dimensions:
\begin{equation}
\varsigma_{n,m} = \frac{1}{Z_{n,m}^{c}} \sum_{z \in \mathcal{Z}_{n,m}^{>0}} \min \left( \frac{\sum_{i \in \operatorname{Mem}(\mathcal{A}_m)} a_{i,m}^{(z)}}{\hat{l}_{n,m}^{(z)}}, 1 \right)
\label{eq:completion_degree}
\end{equation}

where $\mathcal{Z}_{n,m}^{>0} = \{z \mid \hat{l}_{n,m}^{(z)} > 0\}$ is the set of resource dimensions whose expected demand is positive, and $Z_{n,m}^{c} = |\mathcal{Z}_{n,m}^{>0}|$ is the number of effective dimensions, with the satisfaction of each dimension capped at 1.

The individual utility of robot $n$ for participating in coalition $\mathcal{A}_m$ is defined as its accrued revenue minus its total execution cost. Let $\mathcal{V}_m \triangleq \mathcal{V}_{k^*}$ denote the true value of task $m$, where $k^*$ is the true task type and $\mathcal{V}_{k^*}$ is the revenue value associated with that type. Under a proportional allocation rule, the realized utility is:
\begin{equation}
u_n^m(\mathcal{A}_m) = \frac{\|A_m^{(n)}\|_1}{\sum\limits_{i \in \operatorname{Mem}(\mathcal{A}_m)} \|A_m^{(i)}\|_1} \Big( \mathcal{V}_{m} \cdot \varsigma_{n,m} \Big) - E_{n,m}^{Cost}
\label{eq:true_utility}
\end{equation}
where the $\frac{\|A_m^{(n)}\|_1}{\sum_{i \in \operatorname{Mem}(\mathcal{A}_m)} \|A_m^{(i)}\|_1}$ is the relative resource contribution of robot $n$ within the coalition.

However, since the true task type $k^*$ is unknown, each robot $n$ replaces $\mathcal{V}_m$ with the conditional expectation under its local belief, $\widehat{\mathcal{V}}_{n,m} = \mathbb{E}_n[\mathcal{V}_{k^*}] = \sum_{k \in \mathcal{K}} b_{n,k}^m \cdot \mathcal{V}_k$, and evaluates the expected individual utility:
\begin{equation}
u_n^m(\mathcal{A}_m) = \frac{\|A_m^{(n)}\|_1}{\sum\limits_{i \in \operatorname{Mem}(\mathcal{A}_m)} \|A_m^{(i)}\|_1} \Big( \widehat{\mathcal{V}}_{n,m} \cdot \varsigma_{n,m} \Big) - E_{n,m}^{Cost}
\label{eq:expected_utility}
\end{equation}
As robots hold heterogeneous beliefs $\mathcal{B}_n^m$, the resulting $u_n^m$ may differ across coalition members. All subsequent decisions are based on this expected utility.

Furthermore, the comprehensive execution cost, $E_{n,m}^{Cost}$, is decomposed as:
\begin{equation}
E_{n,m}^{Cost} = E_{n,m}^{move} + E_{n,m}^{wait} + E_{n,m}^{exec}
\label{eq:total_cost}
\end{equation}
where $E_{n,m}^{move}$, $E_{n,m}^{wait}$, and $E_{n,m}^{exec}$ correspond to the movement cost, collaborative waiting cost, and execution cost, respectively. These components are detailed as follows.
\subsubsection{Movement Cost}

Because each robot may participate in multiple tasks, the execution order affects the arrival times and the final total utility. To standardize the task-execution process, each robot arranges its task sequence in descending order of relative task importance, which in turn determines the arrival time at each task. The energy consumed when robot $n$ moves from the previous task $m'$ to the current task $m$ is denoted by $E_{n,m}^{move}$:
\begin{equation}
E_{n,m}^{move} = \varpi \cdot Dist(m', m)
\label{eq:cost_move}
\end{equation}
where $Dist(m', m)$ denotes the spatial distance between tasks $m'$ and $m$, and $\varpi$ is the energy-consumption coefficient per unit distance.
\subsubsection{Cooperative Waiting Cost}
During multi-robot cooperative execution, coalition members must remain synchronized. The waiting cost therefore consists of the waiting time before synchronization and the waiting time after synchronization:
\begin{equation}
E_{n,m}^{wait} = \gamma \cdot \left( T_{n,m}^{\mathrm{pre}} + T_{n,m}^{\mathrm{post}} \right)
\label{eq:cost_wait}
\end{equation}
Here, the waiting time of robot $n$ before synchronized execution starts at task $m$ is defined as $T_{n,m}^{\mathrm{pre}} = ST_m - AT_{n,m}$, namely, the time from arrival at the task location to the common start of coalition execution. The synchronized start time of the coalition is determined by the latest-arriving member, namely $ST_m = \max_{i \in \operatorname{Mem}(\mathcal{A}_m)} AT_{i,m}$. Further, let $T_{n,m}^{\mathrm{exec}}$ denote the duration of the portion of task $m$ executed by robot $n$. Then the time when robot $n$ completes its own assigned part is $FT_{n,m} = ST_m + T_{n,m}^{\mathrm{exec}}$, and the time when the coalition finishes as a whole is $FT_m = \max_{i \in \operatorname{Mem}(\mathcal{A}_m)} FT_{i,m}$. Therefore, the waiting time after synchronization can be written as $T_{n,m}^{\mathrm{post}} = FT_m - FT_{n,m}$, namely, the waiting time from when robot $n$ finishes its own part to when the coalition completes as a whole.

\subsubsection{Execution Cost}
During task execution, the robot completes its assigned portion through continuous information collection and data transmission. The execution cost is therefore proportional to the actual execution duration $T_{n,m}^{\mathrm{exec}}$, namely $E_{n,m}^{exec} = \beta \cdot T_{n,m}^{\mathrm{exec}}$, where $\beta$ denotes the energy-consumption coefficient per unit time.

Fig.~\ref{fig:task_schedule_timing} presents the timeline of coalition execution for task $m$. The blue, pink, and green regions denote the movement, waiting, and execution stages, respectively, and the figure highlights the key quantities $AT_{n,m}$, $ST_m$, $FT_{n,m}$, $FT_m$, $T_{n,m}^{\mathrm{pre}}$, $T_{n,m}^{\mathrm{exec}}$, and $T_{n,m}^{\mathrm{post}}$.

\begin{figure}[t]
\centering
\includegraphics[width=1\columnwidth]{sche.pdf}
\caption{Timing diagram of task $m$ execution.}
\label{fig:task_schedule_timing}
\end{figure}

\subsection{Problem Formulation}
Our objective is to determine the optimal overlapping coalition structure $\mathcal{SC}^*$ that maximizes the global total utility, i.e., the sum of all individual robot utilities. The optimization problem is formulated as

\begin{align}
    \max_{\mathcal{SC}} \quad & U = \sum_{m \in \mathcal{T}} \sum_{n \in \operatorname{Mem}(\mathcal{A}_m)} u_n^m (\mathcal{A}_m) \label{eq:objective} \\
    \text{s.t.} \quad
    & \sum_{m \in \mathcal{T}} E_{n,m}^{Cost} \le E_n^{max}, \quad \forall n \in \mathcal{R} \label{eq:constraint_energy} \\
    & 0 \le a_{n,m}^{(z)} \le \eta_n^{(z)}, \quad \forall n \in \mathcal{R}, \forall m \in \mathcal{T}, \forall z \in \mathcal{Z} \label{eq:constraint_resource}
\end{align}

where constraint (\ref{eq:constraint_energy}) represents the energy constraint, which requires that the total energy consumed by each agent when executing its assigned task sequence does not exceed its maximum available energy. Constraint (\ref{eq:constraint_resource}) ensures that the actual capability allocated by a robot to a task does not exceed its maximum service capability.

Since the overlapping coalition formation problem formulated above is NP-hard, obtaining the exact optimal coalition structure is computationally challenging. To address this issue, this paper develops a low-complexity heuristic OCF method to efficiently construct near-optimal coalition structures. 
\section{GAME-BASED OVERLAPPING COALITION FORMATION}
\subsection{Overlapping Coalition Formation Game Model}
Overlapping coalition formation (OCF) games allow agents to join multiple coalitions simultaneously and allocate resources or capabilities across different coalitions. This differs from conventional coalition formation games (CFGs), in which a single agent can belong to only one coalition at a time. In the considered problem, heterogeneous robots possess various reusable capabilities and can simultaneously serve multiple task coalitions. Accordingly, this paper models the problem as an OCF game and maximizes the overall system utility through distributed task allocation based on robot preference relations.

\textbf{Definition 1 (Overlapping Coalition Formation Game Model):}
An overlapping coalition formation (OCF) game for multi-robot collaborative task allocation is defined as the tuple
\begin{align}
\mathcal{G}
=
\left\langle
\mathcal{R},
\mathcal{T},
\{\mathcal{X}_n\}_{n\in\mathcal{R}},
\mathcal{SC},
\mathcal{U}
\right\rangle,
\end{align}
where the elements are defined as follows:

\begin{itemize}
    \item \textbf{Players $\mathcal{R}$:}
    Each robot $n\in\mathcal{R}$ acts as a rational player with autonomous decision-making capability.

    \item \textbf{Task Set $\mathcal{T}$:}
    Each task $m\in\mathcal{T}$ is served by a coalition of robots that jointly invest capabilities into it.

    \item \textbf{Strategy Space $\mathcal{X}_n$:}
    The strategy of robot $n$ is its capability-allocation plan across all tasks, denoted by
    \begin{align}
    X_n=\bigl[A_1^{(n)},\dots,A_M^{(n)}\bigr],
    \end{align}
    where $A_m^{(n)}$ is the capability-investment vector of robot $n$ for task $m$. The feasible strategy space of robot $n$ is denoted by $\mathcal{X}_n$, and each feasible strategy must satisfy the capability upper bounds, time budget, and energy constraints.



    \item \textbf{Overlapping Coalition Structure $\mathcal{SC}$:}
An overlapping coalition structure is defined by 
\begin{align}
\mathcal{SC}=\{\mathcal{A}_1,\dots,\mathcal{A}_M\},
\end{align}
where $\mathcal{A}_m$ denotes the capability resource vector allocated to task $m$ by the robots performing task $m$. Since robots can simultaneously allocate different capability dimensions to multiple tasks, the member sets of different coalitions may overlap, i.e.,
\begin{align}
\operatorname{Mem}(\mathcal{A}_m)\cap\operatorname{Mem}(\mathcal{A}_{m'})\neq\emptyset,\quad \forall m\neq m'.
\end{align}

\item \textbf{Coalition Utility Function $\mathcal{U}$:} The coalition utility mapping $\mathcal{U}$ characterizes the overall return generated by each task coalition under a given coalition structure. For any task $m\in\mathcal{T}$, the utility of coalition $\mathcal{A}_m$ is defined as
\begin{align}
U_m(\mathcal{A}_m)=\mathcal{U}(\mathcal{A}_m),
\end{align}
where $U_m(\mathcal{A}_m)$ is the sum of the local utilities of all participating robots, i.e.,
\begin{align}
U_m(\mathcal{A}_m)=\sum_{n\in \operatorname{Mem}(\mathcal{A}_m)} u_n^m(\mathcal{A}_m),
\end{align}
and $u_n^m(\mathcal{A}_m)$ is defined in~(\ref{eq:expected_utility}).

\end{itemize}

\subsection{Preference Relations and Switching Criteria}


In OCF games, the switching criterion determines whether a robot adjusts its coalition participation and reallocates its resources among different coalitions. Common switching criteria include the Pareto criterion and the selfish criterion. However, the Pareto criterion is often too restrictive, since any switch must not reduce the utility of any other robot. In contrast, the selfish criterion considers only individual utility improvement and may lead the system to local optima.

To balance individual rationality and global social utility, this paper adopts the holistic altruistic preference relation as the switching criterion in the OCF game model.

\subsubsection{Preference Relation}

In OCF games, robots establish preference relations by comparing the total net utility under different overlapping coalition structures. Therefore, a robot's coalition selection is essentially guided by its preference for maximizing the total net utility. We first introduce the preference relation as follows.

\textbf{Definition 2 (Preference Relation):}
For any robot $n\in\mathcal{R}$, let $\mathcal{SC}^P$ and $\mathcal{SC}^Q$ be two feasible overlapping coalition structures. If
\begin{align}
U_n(\mathcal{SC}^Q)>U_n(\mathcal{SC}^P),
\end{align}
then robot $n$ is said to prefer coalition structure $\mathcal{SC}^Q$ to $\mathcal{SC}^P$, denoted by
\begin{align}
\mathcal{SC}^Q \succ_n \mathcal{SC}^P.
\end{align}
If
\begin{align}
U_n(\mathcal{SC}^Q)\ge U_n(\mathcal{SC}^P),
\end{align}
then it is denoted by
\begin{align}
\mathcal{SC}^Q \succeq_n \mathcal{SC}^P.
\end{align}

\textbf{Definition 3 (Holistic Altruistic Preference)\cite{fang2026game}:}
For any robot $n \in \mathcal{R}$ and any two feasible overlapping coalition structures $\mathcal{SC}^P$ and $\mathcal{SC}^Q$ that differ only in robot $n$'s capability allocation strategy, robot $n$ is said to holistically prefer $\mathcal{SC}^Q$ over $\mathcal{SC}^P$, denoted by $\mathcal{SC}^Q \succ_n \mathcal{SC}^P$, if and only if
\begin{equation}
\mathcal{SC}^Q \succ_n \mathcal{SC}^P \iff
\left\{
\begin{aligned}
&U_n(\mathcal{SC}^Q) > U_n(\mathcal{SC}^P), \\
&U_n(\mathcal{SC}^Q) - U_n(\mathcal{SC}^P) \\
&\quad > \sum_{o \in \mathcal R\setminus\{n\}}
\bigl(U_o(\mathcal{SC}^P)-U_o(\mathcal{SC}^Q)\bigr).
\end{aligned}
\right.
\label{eq:altruistic_preference}
\end{equation}

Here, $U_n(\mathcal{SC})$ denotes the total net utility of robot $n$ under the coalition structure $\mathcal{SC}$, obtained by summing its utilities over all task coalitions:
\begin{equation}
U_n(\mathcal{SC}) = \sum_{m \in \mathcal{T}} u_n^m(\mathcal{A}_m).
\label{eq:individual_total_utility_SC}
\end{equation}

This criterion requires that a switch not only improve the initiator's utility, but also yield a strict net gain for the overall system. Rearranging the second inequality gives
\begin{align}
&U_n(\mathcal{SC}^Q) + \sum_{o \in \mathcal{R} \setminus \{n\}} U_o(\mathcal{SC}^Q) \nonumber \\
&\quad > U_n(\mathcal{SC}^P) + \sum_{o \in \mathcal{R} \setminus \{n\}} U_o(\mathcal{SC}^P), \nonumber \\
&\Rightarrow \sum_{r \in \mathcal{R}} U_r(\mathcal{SC}^Q) > \sum_{r \in \mathcal{R}} U_r(\mathcal{SC}^P),
\end{align}
which shows that every accepted switch strictly increases the global social utility. Therefore, this criterion relaxes the Pareto criterion while preserving individual rationality.

\subsubsection{Switching Operation and Strategy}

\textbf{Definition 4 (o-profitable Deviation):}
For the current overlapping coalition structure $\mathcal{SC}^P$, a unilateral capability-reallocation operation by robot $n$ that generates a new coalition structure $\mathcal{SC}^Q$ is called an \textbf{o-profitable deviation} if and only if it satisfies the energy and capability constraints and the total system utility does not decrease:
\begin{align}
\sum_{r \in \mathcal{R}} U_r(\mathcal{SC}^Q) \geq \sum_{r \in \mathcal{R}} U_r(\mathcal{SC}^P).
\label{eq:o_profitable}
\end{align}

\textbf{Definition 5 (Switching Operation):}
For a given overlapping coalition structure $\mathcal{SC}^P$, a switching operation occurs when robot $n$ transfers its $z$th capability component $\eta_n^{(z)}$ from coalition $\mathcal{A}_i^P$ to coalition $\mathcal{A}_j^P$. The updated coalitions are
\begin{equation}
\mathcal{A}_i^Q=\mathcal{A}_i^P\setminus\{\eta_n^{(z)}\},\qquad
\mathcal{A}_j^Q=\mathcal{A}_j^P\cup\{\eta_n^{(z)}\},
\end{equation}
and the resulting coalition structure is
\begin{equation}
\mathcal{SC}^Q=\mathcal{SC}^P\setminus\{\mathcal{A}_i^P,\mathcal{A}_j^P\}\cup\{\mathcal{A}_i^Q,\mathcal{A}_j^Q\}.
\label{eq:switch_structure}
\end{equation}
The switching operation is o-profitable if it satisfies Definition~4.

\subsubsection{Stable Overlapping Coalition Structure}

\textbf{Definition 6 (o-stable Overlapping Coalition Structure):}
For the current overlapping coalition structure
\begin{align}
\mathcal{SC}^P=\{\mathcal{A}_1^P,\dots,\mathcal{A}_M^P\},
\end{align}
if there exists no unilateral o-profitable deviation that satisfies the constraints for any robot $n\in\mathcal{R}$, then $\mathcal{SC}^P$ is called an o-stable overlapping coalition structure.

\section{ALGORITHM FOR OVERLAPPING COALITION FORMATION}
Based on the OCF game model described above, this paper proposes the BHTA algorithm to address the problem in a distributed framework. The proposed method consists of three main phases: heuristic coalition search, task execution and cooperative observation, and information aggregation with belief updating. The detailed procedure is presented in Algorithm~\ref{alg:main_framework}, and Fig.~\ref{fig:bhta_framework} illustrates the overall closed-loop framework.

1) \textbf{Heuristic coalition-search phase:} At round $t$, under the current local beliefs, the system invokes Algorithm~2 to perform heuristic coalition search based on tabu-list simulated annealing and obtain a high-quality overlapping coalition structure for the current round.

2) \textbf{Task execution and cooperative observation phase:} The robot swarm executes the allocated tasks according to the current coalition structure, while collecting local observation information during task execution.

3) \textbf{Information aggregation and belief-update phase:} As shown in Algorithm~3, the robot swarm aggregates the newly acquired observations over the network and updates the local beliefs based on the belief updating mechanism, thereby providing refined estimates of task value and task demand for the next round of coalition search.

The above procedure is repeated over successive rounds until the game reaches termination, after which the uncertainty is resolved, the true task demands become clear, and a stable overlapping coalition structure is obtained and executed by the robot swarm.
\begin{algorithm}[!t]
\caption{BHTA dynamic game framework with cooperative observation and belief updating}
\label{alg:main_framework}
\begin{algorithmic}[1]
\REQUIRE Robot set $\mathcal{R}$, task set $\mathcal{T}$, task-type set $\mathcal{K}$, maximum number of game rounds $T_{max}$, initial temperature $T_0$, cooling coefficient $\alpha$, minimum temperature threshold $T_{min}$
\ENSURE Evolved global overlapping coalition structure $\mathcal{SC}^*$ and final local belief set $\{\mathcal{B}_n^m(T_{max})\}_{n,m}$

\STATE \textbf{Initialization:}
\STATE Initialize the game round as $t \leftarrow 1$, the system temperature as $Temp \leftarrow T_0$, and the initial empty coalition structure as $\mathcal{SC}(0) \leftarrow \emptyset$;
\FOR{$\forall n \in \mathcal{R}$ \AND $\forall m \in \mathcal{T}$}
    \STATE Initialize the local belief vector $\mathcal{B}_n^m(0)$ using the uniform initialization induced by a uniform Dirichlet prior when no prior knowledge is available;
    \STATE Initialize the local observation-count matrix $O_n(0) \leftarrow \mathbf{0}_{M \times K}$;
\ENDFOR

\WHILE{$t \le T_{max}$}
    \STATE \COMMENT{\textbf{Phase 1: overlapping coalition formation game}}
    \STATE Inherit the coalition structure from the previous round $\mathcal{SC}(t-1)$ as a warm start, and call \textbf{Algorithm 2} with the current temperature $Temp$ and local belief set $\{\mathcal{B}_n^m(t)\}_{n,m}$ to solve for the optimal coalition structure $\mathcal{SC}(t)$ of the current round;

    \STATE \COMMENT{\textbf{Phase 2: physical execution and cooperative observation}}
    \STATE Deploy the robot swarm to each task node according to $\mathcal{SC}(t)$, invest resources, and execute local observation to acquire new task-type labels;

    \STATE \COMMENT{\textbf{Phase 3: information fusion and belief updating}}
    \STATE Call \textbf{Algorithm 3} to fuse the observation labels across the network and update the local observation statistics and posterior local beliefs $\{\mathcal{B}_n^m(t+1)\}_{n,m}$ of all robots;

    \STATE Update the temperature as $Temp \leftarrow \max(\alpha \cdot Temp, T_{min})$ and the round index as $t \leftarrow t + 1$;
\ENDWHILE

\STATE \textbf{Output:} Upon reaching the maximum number of game rounds, output the final coalition structure $\mathcal{SC}(T_{max})$ and the local belief set $\{\mathcal{B}_n^m(T_{max})\}_{n,m}$
\end{algorithmic}
\end{algorithm}

\subsection{Heuristic Coalition Formation Formation Method}

Existing studies on overlapping coalition formation usually use unguided best-response or random-search methods, and further develop approaches such as Preference Gravity-Guided Tabu Search (PGG-TS) to improve search guidance and avoid revisiting explored solutions. However, in the scenario considered in this paper, these methods may still suffer from slow convergence and local stagnation. To address this issue, this paper introduces simulated annealing into the distributed coalition-formation process and combines it with a tabu-memory mechanism to design a heuristic coalition-formation strategy. In each game round, robots estimate task values and demands based on local beliefs, and generate resource-reallocation solutions according to heuristic scores and Boltzmann probabilities. At the high-temperature stage, suboptimal moves are accepted with a certain probability to improve global search. At the low-temperature stage, the search gradually shifts to local search around promising coalition structures, while the tabu list prevents cyclic revisits and repeated exploration. Through the coordination of these mechanisms, the proposed method can obtain high-quality overlapping coalition structures more efficiently while maintaining computational feasibility. The detailed steps are given below.

\subsubsection{Heuristic Feature Extraction and Composite Scoring}
To support capability-allocation decisions, each robot evaluates task attractiveness based on its available resources. Since true task demands are unknown and direct optimization is difficult, the algorithm uses a heuristic composite score that jointly considers unmet task demand and robot supply, thereby guiding probabilistic task selection toward suitable coalitions.

At the $k$th iteration, the overlapping coalition structure is denoted by  $\mathcal{SC}^{(k)} = \{ \mathcal{A}_1^{(k)}, \ldots, \mathcal{A}_M^{(k)} \} $, where the capability vector allocated by robot $n$ to task $m$ is denoted by $A_m^{(n)}(k) = \left[ a_{n,m}^{(1)}(k), \ldots, a_{n,m}^{(z)}(k), \ldots, a_{n,m}^{(Z)}(k) \right] $.
For individual decision-making, since the true task requirements are unknown, robot $n$ estimates the expected demand of task $m$ as $\widehat{\mathcal{L}}_{n,m}$ based on its local belief $\mathcal{B}_n^m$. Accordingly, the estimated capability gap of task $m$ in dimension $z$ is defined as
\begin{align}
Gap_{n,m}^{(z)}(k) = \max \Big( \hat{l}_{n,m}^{(z)} - \sum_{i \in \operatorname{Mem}(\mathcal{A}_m^{(k)})} a_{i,m}^{(z)}(k), 0 \Big).
\end{align}

To make physical indicators with different units comparable in decision making, the strategy first normalizes the four key features that determine task attractiveness. For brevity, the iteration index $(k)$ is explicitly retained only for quantities that vary with the search state:

Specifically, the task priority $\mathit{Pri}_m$ measures the urgency of task $m$ and is defined as $\mathit{Pri}_m = \rho_m / \rho_{max}$. The resource shortage $\mathit{Dem}_{n,m}^{(z)}(k)$ measures, from robot $n$'s perspective, the estimated capability gap of task $m$ in dimension $z$, and is defined as $\mathit{Dem}_{n,m}^{(z)}(k) = Gap_{n,m}^{(z)}(k) / Gap_{n,\max}^{(z)}(k)$, where $Gap_{n,\max}^{(z)}(k) = \max_{j \in \mathcal{T}} Gap_{n,j}^{(z)}(k)$. The robot supply $\mathit{Sup}_n^{(z)}$ measures the available capability of robot $n$ in dimension $z$ and is defined as $\mathit{Sup}_n^{(z)} = \eta_n^{(z)} / \eta_{max}^{(z)}$. The spatial distance $\mathit{Dis}_{n,m}(k)$ measures the movement cost for robot $n$ from its current location or previous task location to task $m$, and is defined as $\mathit{Dis}_{n,m}(k) = Dist_{n,m}(k) / Dist_{max}$.

Based on the above features, the heuristic composite score of robot $n$ for investing resource $z$ into task $m$ at the $k$th iteration, denoted by $S_{n,m}^{(z)}(k)$, is defined as

\begin{equation}
S_{n,m}^{(z)}(k) = \frac{\mathit{Pri}_m^2 \cdot \mathit{Dem}_{n,m}^{(z)}(k) \cdot \mathit{Sup}_n^{(z)}}{\mathit{Dis}_{n,m}(k) + \epsilon}
\label{eq:heuristic_score}
\end{equation}

where $\epsilon$ is a small positive number introduced to prevent division by zero. This scoring mechanism assigns higher weight to high-priority tasks and favors tasks with a larger estimated gap, stronger dispatchable capability, and lower movement cost.

\subsubsection{Boltzmann Probability Selection}

The composite scores are converted into task-selection probabilities through the Boltzmann distribution. Specifically, the probability that robot $n$ allocates its $z$th capability to task $m$ at iteration $k$ is given by
\begin{equation}
p_{n,m}^{(z)}(k) =
\begin{cases}
\frac{\exp\left( S_{n,m}^{(z)}(k) / T(k) \right)}{\sum\limits_{j \in \Omega_n^{(z)}(k)} \exp\left( S_{n,j}^{(z)}(k) / T(k) \right)}, & \text{if } m \in \Omega_n^{(z)}(k), \\[6pt]
0, & \text{otherwise},
\end{cases}
\label{eq:boltzmann_prob}
\end{equation}
where $T(k)$ denotes the system temperature at iteration $k$, and
\[
\Omega_n^{(z)}(k) = \{j \in \mathcal{T} \mid S_{n,j}^{(z)}(k) > \epsilon\}
\]
denotes the feasible candidate-task set after excluding tasks for which robot $n$ has no relevant capability or whose estimated gap has already been filled. The resulting selection-probability vector for robot $n$'s $z$th capability over all tasks is denoted by
\begin{align}
    P_n^{(z)}(k) = \left[ p_{n,1}^{(z)}(k), \ldots, p_{n,M}^{(z)}(k) \right].
    \label{xiangliang}
\end{align}

At high temperature, the probability distribution remains relatively broad, which supports global exploration. As the temperature decreases, the probability mass gradually concentrates on high-score tasks, leading the search toward local convergence. To improve numerical stability, the implementation uses a shifted softmax by subtracting the maximum score before exponentiation.

After obtaining the resource-task selection probabilities, the algorithm performs one round of internal search by combining random resource withdrawal, candidate coalition-structure generation, Metropolis-based acceptance, and tabu-memory mechanisms. The search objective is
\begin{equation}
GSU(\mathcal{SC}) \triangleq \sum_{r \in \mathcal{R}} U_r(\mathcal{SC}).
\end{equation}
This process iteratively improves the coalition structure under the current local beliefs and outputs the stable coalition structure found in the current round, as summarized in Algorithm~2.

% \begin{algorithm}[!t]
% \caption{Tabu-Enhanced Simulated Annealing Heuristic for Distributed Overlapping Coalition Formation}
% \label{alg:bhta_internal_search}
% \begin{algorithmic}[1]
% \REQUIRE $\mathcal{R}$, $\mathcal{T}$, $T_{max}$, $T_{min}$, $\alpha$, $p_{leave}$, $L_{tabu}$, $K_{max}^{stable}$
% \STATE \textbf{Initialization:} $k = 1$, $k_{stable} = 0$, $T(1) = T_{max}$;
% \STATE Initialize $\mathit{Tabu}_n = \emptyset$ for each robot $n \in \mathcal{R}$;
% \STATE Generate initial $\mathcal{SC}^{(1)}$; $\mathcal{SC}^* = \mathcal{SC}^{(1)}$, $U^* = GSU(\mathcal{SC}^{(1)})$;

% \WHILE{$k_{stable} \le K_{max}^{stable}$ \AND $T(k) \ge T_{min}$}
%     \STATE $\mathcal{SC} = \mathcal{SC}^{(k)}$; $U = GSU(\mathcal{SC})$;

%     \FOR{each robot $n \in \mathcal{R}$}
%         \STATE Robot $n$ withdraws part of its current coalition participation and releases the corresponding allocated capability resources $\eta_n^{(z)}$ with probability $p_{leave}$;
%         \STATE Robot $n$ computes the selection probability vector $P_n^{(z)}(k)$ according to~(\ref{xiangliang});
%         \STATE Robot $n$ executes the switching operation on $\mathcal{SC}$ according to the computed probability vector, yielding the candidate coalition structure $\mathcal{SC}'$;
%         \STATE $U' = GSU(\mathcal{SC}')$;

%         \IF{$\mathcal{SC}' \in \mathit{Tabu}_n$}
%             \IF{$U' > U^*$}
%                 \STATE $(\mathcal{SC}, U, \mathcal{SC}^*, U^*) \leftarrow (\mathcal{SC}', U', \mathcal{SC}', U')$;
%             \ENDIF
%         \ELSE
%             \IF{$\mathcal{SC}' \succ_n \mathcal{SC}$}
%                 \STATE $(\mathcal{SC}, U) \leftarrow (\mathcal{SC}', U')$;
%                 \STATE Update $\mathit{Tabu}_n$ based on $\mathcal{SC}$;
%                 \IF{$U' > U^*$}
%                     \STATE $(\mathcal{SC}^*, U^*) \leftarrow (\mathcal{SC}', U')$;
%                 \ENDIF
%             \ELSE
%                 \STATE $\Delta E = U' - U$;
%                 \IF{$\mathrm{rand}() \le \exp(\Delta E / T(k))$}
%                     \STATE $(\mathcal{SC}, U) \leftarrow (\mathcal{SC}', U')$;
%                     \STATE Update $\mathit{Tabu}_n$ based on $\mathcal{SC}$;
%                     \IF{$U' > U^*$}
%                         \STATE $(\mathcal{SC}^*, U^*) \leftarrow (\mathcal{SC}', U')$;
%                     \ENDIF
%                 \ENDIF
%             \ENDIF
%         \ENDIF

%         \STATE Pass the updated $\mathcal{SC}$, $\mathcal{SC}^*$, and $U^*$ to the next robot;
%     \ENDFOR

%     \STATE $\mathcal{SC}^{(k+1)} = \mathcal{SC}$;
%     \IF{$\mathcal{SC}^{(k+1)} = \mathcal{SC}^{(k)}$}
%         \STATE $k_{stable} = k_{stable} + 1$;
%     \ELSE
%         \STATE $k_{stable} = 0$;
%     \ENDIF

%     \STATE $T(k+1) = \alpha \cdot T(k)$; $k = k + 1$;
% \ENDWHILE

% \STATE \textbf{Output:} stable coalition structure $\mathcal{SC}^*$.
% \end{algorithmic}
% \end{algorithm}


\begin{algorithm}[!t]
\caption{Tabu-Enhanced Simulated Annealing Heuristic for Distributed Overlapping Coalition Formation}
\label{alg:bhta_internal_search}
\begin{algorithmic}[1]
\REQUIRE $\mathcal{R}$, $\mathcal{T}$, $T_{max}$, $T_{min}$, $\alpha$, $p_{leave}$, $L_{tabu}$, $K_{max}^{stable}$
\STATE \textbf{Initialization:} $k = 1$, $k_{stable} = 0$, $T(1) = T_{max}$, $\mathit{Tabu}_n = \emptyset~(\forall n \in \mathcal{R})$;
\STATE Generate $\mathcal{SC}^{(1)}$ and set $(\mathcal{SC}^*, U^*) = (\mathcal{SC}^{(1)}, GSU(\mathcal{SC}^{(1)}))$;

\WHILE{$k_{stable} \le K_{max}^{stable}$ \AND $T(k) \ge T_{min}$}
    \STATE $\mathcal{SC} = \mathcal{SC}^{(k)}$; $U = GSU(\mathcal{SC})$;

    \FOR{each robot $n \in \mathcal{R}$}
\STATE Robot $n$ withdraws part of its current coalition participation and releases the corresponding allocated capability resources $\eta_n^{(z)}$ with probability $p_{leave}$;
\STATE Robot $n$ computes the selection probability vector $P_n^{(z)}(k)$ according to~(\ref{xiangliang});
\STATE Robot $n$ executes the switching operation on $\mathcal{SC}$ according to the computed probability vector, yielding the candidate coalition structure $\mathcal{SC}'$;
\STATE Robot $n$ applies Algorithm~\ref{alg:acceptance_rule} to decide whether to accept $\mathcal{SC}'$ and update $\mathcal{SC}$, $\mathcal{SC}^*$, and $U^*$;

        \STATE Pass the updated $\mathcal{SC}$, $\mathcal{SC}^*$, and $U^*$ to the next robot;
    \ENDFOR

    \STATE $\mathcal{SC}^{(k+1)} = \mathcal{SC}$;
    \IF{$\mathcal{SC}^{(k+1)} = \mathcal{SC}^{(k)}$}
        \STATE $k_{stable} = k_{stable} + 1$;
    \ELSE
        \STATE $k_{stable} = 0$;
    \ENDIF

    \STATE $T(k+1) = \alpha \cdot T(k)$; $k = k + 1$;
\ENDWHILE

\STATE \textbf{Output:} stable coalition structure $\mathcal{SC}^*$.
\end{algorithmic}
\end{algorithm}

\begin{algorithm}[!t]
\caption{Tabu and Simulated Annealing-Based Acceptance Rule}
\label{alg:acceptance_rule}
\begin{algorithmic}[1]
\REQUIRE Current coalition structure $\mathcal{SC}$, candidate coalition structure $\mathcal{SC}'$, current utility $U$, best utility $U^*$, tabu list $\mathit{Tabu}_n$, temperature $T(k)$
\ENSURE Updated $\mathcal{SC}$, $U$, $\mathcal{SC}^*$, and $U^*$

\STATE $U' = GSU(\mathcal{SC}')$

\IF{$\mathcal{SC}' \in \mathit{Tabu}_n$}
    \IF{$U' > U^*$}
        \STATE $(\mathcal{SC}, U, \mathcal{SC}^*, U^*) \leftarrow (\mathcal{SC}', U', \mathcal{SC}', U')$
    \ENDIF
\ELSIF{$\mathcal{SC}' \succ_n \mathcal{SC}$}
    \STATE $(\mathcal{SC}, U) \leftarrow (\mathcal{SC}', U')$
    \STATE Update $\mathit{Tabu}_n$ based on $\mathcal{SC}$
    \IF{$U' > U^*$}
        \STATE $(\mathcal{SC}^*, U^*) \leftarrow (\mathcal{SC}', U')$
    \ENDIF
\ELSE
    \STATE $\Delta E = U' - U$
    \IF{$\mathrm{rand}() \le \exp(\Delta E / T(k))$}
        \STATE $(\mathcal{SC}, U) \leftarrow (\mathcal{SC}', U')$
        \STATE Update $\mathit{Tabu}_n$ based on $\mathcal{SC}$
        \IF{$U' > U^*$}
            \STATE $(\mathcal{SC}^*, U^*) \leftarrow (\mathcal{SC}', U')$
        \ENDIF
    \ENDIF
\ENDIF

\STATE \textbf{return} $\mathcal{SC}, U, \mathcal{SC}^*, U^*$
\end{algorithmic}
\end{algorithm}







\subsection{Belief Fusion and Local Belief Updating}

During coalition execution in the current round, the robot swarm performs the assigned tasks while collecting new local observation labels. After task execution is completed, the robots return to the initial base station, where the newly acquired observations are shared and fused through periodic communication to update each robot's local belief over task types. As more observations are incorporated over successive rounds, the robots form more consistent and stable beliefs, which in turn provide refined estimates of task values and task demands for the next round of coalition search. This process corresponds to Phase 3 of Algorithm~1.

Let the probability that robot $n \in \mathcal{R}$ correctly identifies the true task type in a single observation be $d_n \in (0,1)$, which characterizes its observation reliability. Each local observation outputs a task-type label. Owing to environmental disturbance and recognition errors, this label deviates from the true type with probability $1-d_n$. Let $o_{n,m}^{(k)}(t)$ denote the cumulative number of times that robot $n$ has judged task $m \in \mathcal{T}$ as type $k \in \mathcal{K}$ by the end of round $t$.

In round $t+1$, the system fuses the newly obtained observation labels across the network through offline or periodic communication. Let $\delta_{n',m}^{(k)}(t+1)$ denote the number of new observations generated by robot $n'$ in the current round that classify task $m$ as type $k$. Then the globally fused cumulative observation record is updated as

\begin{equation}
    o_{n,m}^{(k)}(t+1) = o_{n,m}^{(k)}(t) + \sum_{n' \in \mathcal{R}} \delta_{n',m}^{(k)}(t+1), \quad \forall m \in \mathcal{T}, \forall k \in \mathcal{K}
\label{eq:observation_update}
\end{equation}

To characterize the uncertainty of the task-type probability vector, this paper uses a uniform Dirichlet prior $\text{Dir}(\alpha_0,\ldots,\alpha_0)$ for the local belief $\mathcal{B}_n^m$, where $\alpha_0 > 0$ is a common pseudo-count parameter. Since each observation outputs a discrete type label, the fused type counts follow multinomial count statistics. Hence, the posterior distribution can be written as

\begin{equation}
    \mathcal{B}_n^m(t+1) \sim \text{Dir} \Big( \alpha_0 + o_{n,m}^{(1)}(t+1), \alpha_0 + o_{n,m}^{(2)}(t+1), \dots, \alpha_0 + o_{n,m}^{(K)}(t+1) \Big)
\label{eq:dirichlet_update}
\end{equation}

Using the expectation property of this Dirichlet posterior, the expected probability that robot $n$ judges task $m$ to belong to type $k$, namely the $k$th element in the belief vector $\mathcal{B}_{n}^m = [b_{n,1}^m, \dots, b_{n,K}^m]$, can be written as

\begin{equation}
    b_{n,k}^m(t+1) = \frac{\alpha_0 + o_{n,m}^{(k)}(t+1)}{\sum_{\kappa \in \mathcal{K}} \left( \alpha_0 + o_{n,m}^{(\kappa)}(t+1) \right)}
\label{eq:belief_expectation}
\end{equation}

In the implementation of this paper, $d_n$ serves to characterize the reliability background of the robot observation labels, whereas the posterior update itself still performs conjugate updating based on the fused type counts. Through this mechanism, the robot swarm can gradually form more consistent cognition of task types as the game evolves, thereby providing more stable demand estimates for subsequent coalition updates.

The above information-fusion and posterior-updating process is summarized in Algorithm~3.
\begin{algorithm}[!t]
\caption{Distributed information fusion and local belief updating}
\label{alg:belief_updating}
\begin{algorithmic}[1]
\REQUIRE Historical local observation-statistics set $\{O_n(t)\}_{n \in \mathcal{R}}$, network-wide increment of newly observed labels $\delta(t+1)$
\ENSURE Updated local observation-statistics set $\{O_n(t+1)\}_{n \in \mathcal{R}}$ and posterior local belief set $\{\mathcal{B}_n^m(t+1)\}_{n,m}$

\STATE \COMMENT{\textbf{Information fusion and posterior belief updating}}
\FOR{each robot $n \in \mathcal{R}$}
    \STATE Obtain the newly generated observation labels $\delta(t+1)$ from the whole network during the current physical-execution stage;
    \FOR{each task $m \in \mathcal{T}$}
        \STATE \textit{\% Fuse the global observation labels}
        \STATE Update the local observation matrix $O_n(t+1)$ by accumulating the network-wide increment according to (\ref{eq:observation_update});

        \STATE \textit{\% Compute the expectation of the Dirichlet posterior}
        \FOR{each type $k \in \mathcal{K}$}
            \STATE Compute the expected probability $b_{n,k}^m(t+1)$ of the task type according to (\ref{eq:belief_expectation});
        \ENDFOR
        \STATE Reassemble the updated local belief vector $\mathcal{B}_n^m(t+1)$;
    \ENDFOR
\ENDFOR

\RETURN $\{O_n(t+1)\}_{n \in \mathcal{R}}, \{\mathcal{B}_n^m(t+1)\}_{n,m}$
\end{algorithmic}
\end{algorithm}

\subsection{Stability and Convergence Analysis}

\noindent\textit{1) Finite-step termination:} For Algorithm~2, if the cooling coefficient satisfies $0 < \alpha < 1$, then the temperature sequence $T(k+1)=\alpha T(k)$ decreases monotonically and falls below the threshold $T_{min}$ within a finite number of steps. Meanwhile, the algorithm also introduces the maximum number of stable iterations $K_{max}^{stable}$ as an additional stopping condition. Therefore, whether termination is triggered by temperature decay or by the absence of state changes over several consecutive iterations, the one-round internal search always terminates in a finite number of steps.

\noindent\textit{2) Relation to the o-stable coalition structure:} Definition 6 has already provided the criterion of an o-stable overlapping coalition structure, namely the nonexistence of any unilateral o-profitable deviation that satisfies the constraints. In one-round search with fixed current local beliefs, Algorithm~2 continuously adjusts the coalition structure through heuristic resource reallocation, the Metropolis acceptance criterion, and tabu memory. When the algorithm terminates, the current coalition structure is output as a locally stable solution under the utility function induced by the current beliefs. It should be emphasized that, because the local beliefs in the outer framework continue to be updated across rounds, BHTA produces a sequence of locally stable coalition structures that evolve with cognition, rather than a single globally optimal fixed point in a static environment.

\noindent\textit{3) Complexity analysis:} Let $K_{SA}$ denote the actual number of iterations of Algorithm~2 in one round of internal search. For each robot in each iteration, the heuristic scoring and Boltzmann probability computation over $M$ tasks and $Z$ capability dimensions have complexity approximately $\mathcal{O}(MZ)$. If the scheduling-feasibility verification and global-social-utility evaluation of a candidate solution are denoted by $C_{sched}$, then the one-round complexity of Algorithm~2 can be written as $\mathcal{O}\!\left(K_{SA} N (MZ + C_{sched})\right)$. For Algorithm~3, network-wide observation fusion and posterior belief updating traverse the three dimensions of robots, tasks, and task types, yielding a one-round complexity of $\mathcal{O}(NMK)$. Therefore, the total complexity of BHTA in a single outer evolution round can be expressed as
\[
\mathcal{O}\!\left(K_{SA} N (MZ + C_{sched}) + NMK\right),
\]
and the total complexity over $T_{max}$ evolution rounds is
\[
\mathcal{O}\!\left(T_{max}\left[K_{SA} N (MZ + C_{sched}) + NMK\right]\right).
\]
If the scheduling-feasibility check is approximated as being of the same order as the task scale, the above complexity can be further approximated as polynomial growth in $N$, $M$, $Z$, and $K$.
\section{SIMULATION RESULTS}

\begin{figure*}[!t]
    \centering
    \includegraphics[width=\textwidth]{varyNM.eps}
    \caption{Performance comparison of four algorithms under varying numbers of tasks and agents.
    (a) Final coalition utility versus the number of tasks $M$ with $N=8$.
    (b) Completion rate versus the number of tasks $M$ with $N=8$.
    (c) Final coalition utility versus the number of agents $N$ with $M=10$.
    (d) Completion rate versus the number of agents $N$ with $M=10$.}
    \label{fig:main}
\end{figure*}

\begin{figure*}[!t]
    \centering
    \includegraphics[width=\textwidth]{single_results.eps}
    \caption{An example of multi-robot overlapping coalition formation under uncertainty.
    (a) Task execution route map.
    (b) Scheduling process of the proposed BHTA algorithm.
    (c) Differences between actual task demand and allocated resources.
    (d) Evolution of coalition utility, total coalition cost, and total completion value.}
    \label{single_results}
\end{figure*}

\begin{figure*}[!t]
    \centering
    \includegraphics[width=\textwidth]{single_viz_resource_allocation_matrix.eps}
    \caption{Resource allocation matrix for the single visualization case.}
    \label{fig:resource_allocation}
\end{figure*}

\begin{figure*}[!t]
    \centering
    \includegraphics[width=\textwidth]{single_viz_belief_expected_task_value.eps}
    \caption{Belief and expected task value visualization for the single visualization case.}
    \label{fig:belief_expected_value}
\end{figure*}

\begin{figure*}[!t]
    \centering
    \includegraphics[width=\textwidth]{ablation_convergence_all_conditions.pdf}
    \caption{Convergence curves for the ablation study under all tested conditions.}
    \label{ablation_convergence}
\end{figure*}

\begin{figure*}[!t]
    \centering
    \includegraphics[width=\textwidth]{ablation_endpoint_scatter_all_conditions.pdf}
    \caption{Endpoint scatter plot for the ablation study under all tested conditions.}
    \label{ablation_scatter}
\end{figure*}

\subsection{Setup}

All simulations were conducted in a two-dimensional multi-robot task-allocation scenario. Unless otherwise specified, $N$ agents and $M$ tasks were randomly placed in a $100\times100\,\mathrm{m}$ rectangular mission region. Each agent carried $Z$ heterogeneous resource dimensions, and each task was associated with a demand vector, task value, and priority level generated from the parameter ranges in Table~\ref{tab:sim_params}.

During execution, agents selected tasks and allocated resources according to their remaining resources, task attributes, and local belief information. The overlapping coalition formation mechanism then coordinated the distributed decisions and produced the final task allocation. The main experiments evaluate BHTA under varying task-agent ratios, a representative single-case visualization, an ablation study on the belief update mechanism, and repeated trials over multiple random seeds.

The simulations were run on a workstation with an Intel(R) Core(TM) i7-14700 CPU at 2.10 GHz, 32.0 GB memory, and MATLAB.
% TODO: Specify the operating system and MATLAB release before submission if these records are available.

\subsection{Baselines, Fairness Settings, and Evaluation Metrics}

\noindent\textit{1) Baselines:} The proposed algorithm BHTA was compared with three baselines: Non-overlapping Coalition Formation (NOCF)~\cite{huo2024self}, Preference Gravity-Guided Tabu Search (PGG-TS)~\cite{qi2022task}, and JSE-OCF adapted from~\cite{shi2023federated}. For JSE-OCF, the joining, switching, and exiting operations were retained and embedded into the same simulation model to form overlapping coalitions through heuristic rules.

\noindent\textit{2) Fairness settings:} To improve comparability, all algorithms were evaluated under the same problem model, simulation framework, initial beliefs, observation budget, stopping conditions, maximum iteration settings, resource-demand estimation rule, and random seeds where applicable. NOCF was constrained such that each agent could execute only one task and used a self-interested preference rule, whereas PGG-TS and JSE-OCF used the same altruistic preference criterion as that in the overlapping-coalition setting.

\noindent\textit{3) Evaluation metrics:} The evaluation focuses on task completion, system utility, and algorithmic robustness. The completion rate (CR) measures the degree to which task demands are satisfied, while the final coalition utility $U$ measures the net system benefit after coalition formation and resource allocation. Convergence curves and endpoint scatter plots are used to examine the evolution and variability of the algorithms across repeated trials.

\begin{table}[t]
\centering
\caption{Main simulation parameters}
\label{tab:sim_params}
\scriptsize
\setlength{\tabcolsep}{5pt}
\renewcommand{\arraystretch}{1.15}
\begin{tabular}{l l}
\toprule
\textbf{Parameter} & \textbf{Value} \\
\midrule
Mission region $\Omega$ & $100\times100\ \mathrm{m}$ \\
Number of outer game rounds $I_{\mathrm{out}}$ & 100 \\
Maximum inner iterations per round $I_{\mathrm{in}}^{\max}$ & 80 \\
Number of observations $N_{\mathrm{obs}}$ & 50 \\
Number of resource dimensions $Z$ & 6 \\
Task value levels $\mathcal{V}_k$ & $[500,1000,2000]$ \\
Upper bounds of task demand $\bar{d}_k$ & $[2,5,8]$ \\
Resource execution times $\tau_z$ & $[50,65,50,60,35,45]\ \mathrm{s}$ \\
Agent speed $v_n$ & $2\ \mathrm{m/s}$ \\
Positive observation probability $p_n^{+}$ & $[0.75,0.95]$ \\
Maximum energy budget $E_n^{max}$ & $[300,350]$ \\
Motion fuel consumption coefficient $\varpi$ & $1\ \mathrm{m}^{-1}$ \\
Waiting fuel consumption coefficient $\gamma$ & 0.5 \\
Execution cost coefficient $\beta$ & 1 \\
Resource capacity range per dimension $\eta_n^{(z)}$ & $[0,4]$ \\
Task priority range $\rho_m$ & $[1,9]$ \\
\bottomrule
\end{tabular}
\end{table}

\begin{table*}[t]
\caption{OVERALL PERFORMANCE COMPARISON OF FOUR ALGORITHMS}
\label{tab:overall_side_by_side_u_cr}
\centering

% Size control
\newcommand{\HeadSize}{\small}      % Header font size
\newcommand{\DataSize}{\scriptsize} % Data font size
\setlength{\tabcolsep}{4.0pt}       % Column spacing
\renewcommand{\arraystretch}{1.10}  % Row height

\makebox[\textwidth][c]{%
\begin{tabular}{@{}ccccc@{\hspace{1.1em}}ccccc@{}}
\toprule
\HeadSize
& \multicolumn{4}{c}{\textbf{Varying number of agents $N$ ($M=10$)}}
&
& \multicolumn{4}{c}{\textbf{Varying number of tasks $M$ ($N=8$)}} \\
\cmidrule(lr){2-5}\cmidrule(lr){7-10}
$N$ & \textbf{NOCF} & \textbf{PGG-TS} & \textbf{JSE-OCF} & \textbf{BHTA
}& $M$ & \textbf{NOCF} & \textbf{PGG-TS} & \textbf{JSE-OCF} & \textbf{BHTA} \\
\midrule

\DataSize
4
& 3028.82(0.209) & 5813.06(0.626) & 5946.66(0.629) & \textbf{6073.62}(\textbf{0.638})
& 8
& 5522.05(0.561) & 7313.49(0.927) & 7320.92(0.927) & \textbf{7752.60}(\textbf{0.947}) \\

6
& 4225.55(0.326) & 7274.63(0.793) & 6989.80(0.758) & \textbf{7624.52}(\textbf{0.808})
& 10
& 5389.47(0.444) & 8080.93(0.858) & 8288.32(0.875) & \textbf{8501.30}(\textbf{0.889}) \\

8
& 5194.94(0.447) & 7822.48(0.874) & 7846.00(0.868) & \textbf{8198.33}(\textbf{0.908})
& 12
& 6025.54(0.407) & 9572.30(0.841) & 9725.80(0.839) & \textbf{10107.74}(\textbf{0.870}) \\

10
& 6078.10(0.563) & 8088.95(0.928) & 8148.88(0.930) & \textbf{8478.81}(\textbf{0.949})
& 14
& 6575.55(0.348) & 11137.44(0.802) & 11056.05(0.797) & \textbf{11595.05}(\textbf{0.828}) \\

12
& 6818.49(0.685) & 8111.99(0.950) & 8011.73(0.954) & \textbf{8656.21}(\textbf{0.977})
& 16
& 6736.19(0.293) & 12188.15(0.752) & 11768.78(0.710) & \textbf{12471.17}(\textbf{0.784}) \\

14
& 7372.59(0.762) & 8172.81(0.972) & 8196.78(0.975) & \textbf{8566.66}(\textbf{0.986})
& 18
& 6515.22(0.260) & 12232.34(0.705) & 11719.03(0.673) & \textbf{12776.79}(\textbf{0.730}) \\

16
& 7863.55(0.828) & 8155.98(0.984) & 8132.76(0.976) & \textbf{8661.85}(\textbf{0.990})
& 20
& 6969.47(0.217) & 13631.63(0.637) & 13016.94(0.602) & \textbf{14366.54}(\textbf{0.660}) \\
\bottomrule
\end{tabular}%
}

\vspace{1.5ex}
\parbox{0.8\textwidth}{\footnotesize
\textit{Note:} Left: varying agents $N$ ($M=10$). Right: varying tasks $M$ ($N=8$). Entries are formatted as $U(\mathrm{CR})$, where $U$ is final coalition utility and $\mathrm{CR}$ is completion rate. Best values are in bold.
}
\end{table*}

\newlength{\SEw}
\setlength{\SEw}{2.4em} % Width of each Start/End subcolumn

\newcolumntype{C}[1]{>{\centering\arraybackslash}m{#1}}

\newcommand{\AllocCell}[1]{%
  \multicolumn{2}{>{\centering\arraybackslash}m{\dimexpr 2\SEw + 2\tabcolsep\relax}}{#1}%
}

\begin{table*}[t]
\caption{Agent--task schedule with resource allocations, carrying capacities, task demands, and completion degrees. Tasks are ordered by priority in descending order.}
\label{tab:agent_task_schedule_with_carry_demand}
\centering
\scriptsize
\setlength{\tabcolsep}{1.9pt}
\renewcommand{\arraystretch}{1.10}

% Column layout: 1 Agent column + 8 tasks x 2 columns (Start/End) = 17 columns
\begin{tabular}{c *{8}{C{\SEw} C{\SEw}}}
\toprule
\multirow{2}{*}{Agent}
& \multicolumn{2}{c}{T4 (P8, V2000)}
& \multicolumn{2}{c}{T5 (P7, V2000)}
& \multicolumn{2}{c}{T3 (P6, V1000)}
& \multicolumn{2}{c}{T6 (P5, V1000)}
& \multicolumn{2}{c}{T2 (P4, V1000)}
& \multicolumn{2}{c}{T8 (P3, V1000)}
& \multicolumn{2}{c}{T7 (P2, V2000)}
& \multicolumn{2}{c}{T1 (P1, V1000)} \\
\cmidrule(lr){2-3}\cmidrule(lr){4-5}\cmidrule(lr){6-7}\cmidrule(lr){8-9}
\cmidrule(lr){10-11}\cmidrule(lr){12-13}\cmidrule(lr){14-15}\cmidrule(lr){16-17}
& Start & End
& Start & End
& Start & End
& Start & End
& Start & End
& Start & End
& Start & End
& Start & End \\
\midrule

% Agent A1
\multirow{2}{*}{\makecell[c]{A1\\$[4,0,3,1,1,3]$}}
& \AllocCell{[4,0,3,0,1,3]}
& \AllocCell{--}
& \AllocCell{[4,0,3,0,0,3]}
& \AllocCell{--}
& \AllocCell{[4,0,3,0,1,3]}
& \AllocCell{--}
& \AllocCell{--}
& \AllocCell{[4,0,3,0,0,3]} \\
& 12.2 & 77.2
& -- & --
& 98.2 & 163.2
& -- & --
& 190.9 & 240.9
& -- & --
& -- & --
& 265.1 & 315.1 \\
\midrule

% Agent A2
\multirow{2}{*}{\makecell[c]{A2\\$[2,4,4,3,0,2]$}}
& \AllocCell{--}
& \AllocCell{[2,4,4,3,0,2]}
& \AllocCell{--}
& \AllocCell{[2,4,4,3,0,2]}
& \AllocCell{--}
& \AllocCell{[2,4,4,3,0,2]}
& \AllocCell{--}
& \AllocCell{--} \\
& -- & --
& 26.9 & 91.9
& -- & --
& 104.0 & 169.0
& -- & --
& 181.7 & 246.7
& -- & --
& -- & -- \\
\midrule

% Agent A3
\multirow{2}{*}{\makecell[c]{A3\\$[0,1,0,2,3,2]$}}
& \AllocCell{--}
& \AllocCell{[0,1,0,0,3,2]}
& \AllocCell{[0,1,0,2,3,0]}
& \AllocCell{--}
& \AllocCell{--}
& \AllocCell{--}
& \AllocCell{[0,1,0,2,3,2]}
& \AllocCell{--} \\
& -- & --
& 26.9 & 91.9
& 98.2 & 163.2
& -- & --
& -- & --
& -- & --
& 187.4 & 252.4
& -- & -- \\
\midrule

% Agent A4
\multirow{2}{*}{\makecell[c]{A4\\$[1,0,0,4,3,0]$}}
& \AllocCell{--}
& \AllocCell{[1,0,0,4,3,0]}
& \AllocCell{--}
& \AllocCell{[1,0,0,4,3,0]}
& \AllocCell{--}
& \AllocCell{[1,0,0,4,3,0]}
& \AllocCell{--}
& \AllocCell{--} \\
& -- & --
& 26.9 & 91.9
& -- & --
& 104.0 & 169.0
& -- & --
& 181.7 & 246.7
& -- & --
& -- & -- \\
\midrule

% Agent A5
\multirow{2}{*}{\makecell[c]{A5\\$[0,2,0,3,2,2]$}}
& \AllocCell{[0,2,0,3,2,2]}
& \AllocCell{--}
& \AllocCell{[0,2,0,3,2,0]}
& \AllocCell{--}
& \AllocCell{--}
& \AllocCell{--}
& \AllocCell{[0,2,0,3,2,2]}
& \AllocCell{--} \\
& 12.2 & 77.2
& -- & --
& 98.2 & 163.2
& -- & --
& -- & --
& -- & --
& 187.4 & 252.4
& -- & -- \\
\midrule

% Demand row
\textbf{Demand}
& \AllocCell{[5,5,6,7,4,5]}
& \AllocCell{[5,5,6,7,4,5]}
& \AllocCell{[3,3,2,5,5,2]}
& \AllocCell{[3,3,2,5,5,2]}
& \AllocCell{[3,3,2,5,5,2]}
& \AllocCell{[3,3,2,5,5,2]}
& \AllocCell{[5,5,6,7,4,5]}
& \AllocCell{[3,3,2,5,5,2]} \\
\midrule

% Allocated row
\textbf{Allocated}
& \AllocCell{[4,2,3,3,3,5]}
& \AllocCell{[3,5,4,7,6,4]}
& \AllocCell{[4,3,3,5,5,3]}
& \AllocCell{[3,4,4,7,3,2]}
& \AllocCell{[4,0,3,0,1,3]}
& \AllocCell{[3,4,4,7,3,2]}
& \AllocCell{[0,3,0,5,5,4]}
& \AllocCell{[4,0,3,0,0,3]} \\
\midrule

% Completion row
\textbf{Completion}
& \AllocCell{64.6\%}
& \AllocCell{84.4\%}
& \AllocCell{100.0\%}
& \AllocCell{93.3\%}
& \AllocCell{53.3\%}
& \AllocCell{93.3\%}
& \AllocCell{51.9\%}
& \AllocCell{50.0\%} \\
\bottomrule
\end{tabular}

\vspace{1ex}
\parbox{0.98\textwidth}{\footnotesize
\textit{Note:} $Q$ denotes the carrying-capacity vector of each agent in the form $[R_1,R_2,R_3,R_4,R_5,R_6]$.
For each agent, the first row gives the allocated resource vector for each executed task, and the second row gives the corresponding start and end times.
\textit{Demand} denotes the task-demand vector.
\textit{Allocated} denotes the total allocated resource vector for each task across all participating agents.}
\end{table*}

\subsection{Overall Comparison}

Table~\ref{tab:overall_side_by_side_u_cr} and Fig.~\ref{fig:main} compare the final coalition utility and CR of BHTA, NOCF, PGG-TS, and JSE-OCF under different values of $N$ and $M$. Across all tested settings, BHTA consistently achieves the highest utility and CR among the four algorithms.

When the number of tasks increases with $N=8$ and $M\in[8,20]$, the available task value increases, and the utility of all algorithms generally rises. At the same time, CR decreases as the fixed robot team becomes more heavily loaded. Under this setting, BHTA improves utility over the best competing baseline by 2.3\%--5.9\%. For example, at $M=20$, BHTA reaches a utility of $14366.54$ and a CR of $0.660$, compared with $13631.63$ and $0.637$ for the best baseline.

When the number of agents increases with $M=10$ and $N\in[4,16]$, the additional resources improve task coverage and increase both utility and CR. BHTA maintains a 2.1\%--6.7\% utility gain over the strongest baseline. The largest relative gain in this group occurs at $N=12$, where BHTA obtains $8656.21$ compared with $8111.99$ for the best baseline. When $N\ge14$, CR becomes close to saturation, but BHTA still improves utility by producing more effective coalition structures.

\subsection{Single-case Visualization}

Fig.~\ref{single_results} illustrates one representative run with five robots and eight tasks under uncertainty. The route map shows the initial spatial distribution of tasks and robots, while the scheduling panel shows how BHTA forms overlapping coalitions during execution. For instance, robots 2 and 3 first cooperate on task 5; after that task is completed, robot 2 moves to task 6, whereas robot 3 proceeds to task 3. Fig.~\ref{single_results}(c) and Fig.~\ref{fig:resource_allocation} further show the mismatch between actual task demands and allocated resources, where positive and negative deviations correspond to over-allocation and under-allocation, respectively.

Table~\ref{tab:agent_task_schedule_with_carry_demand} reports the resulting schedule, carrying capacities, task demands, allocated resources, and completion degrees in descending priority order. The completion row indicates that T3 is fully completed, while T6 and T8 each reach $93.3\%$ and T5 reaches $84.4\%$. Several tasks remain resource-limited, including T4 ($64.6\%$), T2 ($53.3\%$), T7 ($51.9\%$), and T1 ($50.0\%$). Thus, this case illustrates both effective high-value allocation and residual resource gaps under limited agent capacity.

As the game rounds progress in this representative case, the total completed value increases from $5551.59$ to $7919.84$, and the coalition utility increases from $4059.92$ to $6519.07$. Meanwhile, the global cost decreases from $1491.67$ to $1400.77$. Fig.~\ref{fig:belief_expected_value} shows that the task-type beliefs and expected task values become more stable as observations accumulate, which explains the improved allocation quality over the iterative process.

\subsection{Ablation Study}

To isolate the contribution of the belief update mechanism, the complete BHTA is compared with a variant that removes belief updating and relies only on the initial beliefs for task allocation. The number of tasks is fixed at $M=15$, and the number of agents varies over $N\in\{4,8,12,15\}$. Fig.~\ref{ablation_convergence} presents the convergence curves, and Fig.~\ref{ablation_scatter} presents the endpoint distributions of final utility and CR.

\begin{table}[htbp]
\centering
\caption{Ablation Study Results: Utility and Completion Rate across Different Scenarios}
\label{tab:ablation_results_compact}
\footnotesize
\setlength{\tabcolsep}{2.5pt} % Compress column spacing to fit the single-column layout
\renewcommand{\arraystretch}{1} % Slightly tighten row spacing
\begin{tabular}{l c c c c}
\toprule
\multirow{2}{*}{\textbf{Metric}} & \multicolumn{4}{c}{\textbf{Number of Agents ($N$)}} \\
\cmidrule(lr){2-5}
 & \textbf{4} & \textbf{8} & \textbf{12} & \textbf{15} \\ % Simplified to pure numbers to save width
\midrule
$U_{\text{base}}$   & $4505 \pm 1425$ & $6208 \pm 1756$ & $6454 \pm 1974$ & $6372 \pm 2029$ \\
$U_{\text{belief}}$ & $\mathbf{6221 \pm 1669}$ & $\mathbf{9244 \pm 1719}$ & $\mathbf{10515 \pm 1769}$ & $\mathbf{10767 \pm 2049}$ \\
$\Delta U(\%)$      & $+38.1$         & $+48.9$         & $+62.9$          & $+69.0$          \\
\midrule
$CR_{\text{base}}$   & $0.419 \pm 0.084$ & $0.604 \pm 0.100$ & $0.657 \pm 0.102$ & $0.665 \pm 0.107$ \\
$CR_{\text{belief}}$ & $\mathbf{0.464 \pm 0.090}$ & $\mathbf{0.726 \pm 0.078}$ & $\mathbf{0.841 \pm 0.064}$ & $\mathbf{0.877 \pm 0.065}$ \\
$\Delta CR$          & $+0.045$          & $+0.121$          & $+0.184$          & $+0.213$          \\
\bottomrule
\end{tabular}
\vspace{4pt}
\parbox{\columnwidth}{\footnotesize \textbf{Note:} Results are presented as Mean $\pm$ Std over 51 seeds. Best results are highlighted in bold. $U_{\text{base}}$ and $CR_{\text{base}}$ denote the baseline algorithm without belief update, while $U_{\text{belief}}$ and $CR_{\text{belief}}$ represent BHTA. $\Delta U(\%)$ is the relative percentage improvement in utility, and $\Delta CR$ is the absolute difference in CR.}
\end{table}

Table~\ref{tab:ablation_results_compact} quantifies the final performance gap between the two variants. As $N$ increases, the utility improvement from belief updating grows from $38.1\%$ at $N=4$ to $69.0\%$ at $N=15$. The CR gain also increases from $0.045$ to $0.213$. At $N=15$, BHTA achieves an average CR of $0.877$, compared with $0.665$ for the variant without belief updating.

These results indicate that stale or inaccurate initial beliefs can lead the robots to misestimate task types and resource requirements, thereby reducing both completed value and coalition utility. By continuously incorporating observations, BHTA updates the expected task values and demand estimates, which improves resource allocation especially when more agents participate in overlapping coalitions.

\subsection{Convergence and Robustness Discussion}

The convergence curves in Fig.~\ref{ablation_convergence} show that the belief-update version reaches higher utility levels than the no-update variant under the tested conditions. The endpoint scatter plots in Fig.~\ref{ablation_scatter}, together with the mean and standard deviation over 51 seeds in Table~\ref{tab:ablation_results_compact}, further indicate that the performance gain is not limited to a single run. These observations support the use of the belief update mechanism as a recurring component of BHTA, while the statistical scope remains limited to the reported simulation settings.
\section{CONCLUSION}

This paper investigated multi-robot task allocation under uncertain task types and heterogeneous reusable resources, and formulated the problem as an overlapping coalition formation framework. The proposed BHTA integrates local belief estimation, resource-contribution proportional allocation, and a holistic altruistic preference criterion into a dynamic coalition-formation process. Within this framework, robots update task-type beliefs from accumulated observations, estimate task values and resource demands from the current belief state, and search for high-quality overlapping coalition structures using a simulated-annealing-based heuristic with tabu enhancement. This design allows robots to participate in multiple task coalitions while keeping task allocation, resource feasibility, and energy-related execution costs within one utility-driven optimization model.

Simulation results under varying task-agent ratios show that BHTA consistently achieves higher final coalition utility and CR than NOCF, PGG-TS, and JSE-OCF in the tested settings. When the number of tasks increases with $N=8$, BHTA improves utility over the strongest baseline by 2.3\%--5.9\%; when the number of agents increases with $M=10$, the utility gain remains 2.1\%--6.7\%. The representative single-case visualization further illustrates how overlapping coalitions evolve as beliefs and expected task values become more stable. The ablation study over 51 random seeds indicates that the belief update mechanism contributes substantially to both utility and CR, with the utility improvement increasing from 38.1\% to 69.0\% as $N$ grows from 4 to 15. Future work will further examine larger-scale deployments, communication constraints, and real-time replanning under more dynamic task arrivals.

\section{REFERENCES}

\bibliographystyle{IEEEtran}
\bibliography{reference}

\end{document}
 
